\subsection*{7.5 Expectation of the Product of Two Independent Random Variables}

\begin{thmbox}{7.13}
Suppose $X$ and $Y$ are random variables defined on a common probability space $(\Omega,\mathcal{F},P)$. If $X$ and $Y$ are independent and if $E[X]$, $E[Y]$, and $E[XY]$ are all finite, then
\begin{equation}
E[XY]=E[X]E[Y].
\tag{7.10}
\end{equation}
\end{thmbox}

\textit{Proof} We first suppose that $X$ and $Y$ are indicator functions. Suppose $X=1_A$ and $Y=1_B$ for some sets $A$ and $B$. Since $X$ and $Y$ are measurable, the sets $A$ and $B$ are both $\mathcal{F}$-measurable. Moreover, since $X$ and $Y$ are independent, the events $A$ and $B$ are independent. Therefore,
\[
E[1_A1_B] = E[1_{A\cap B}] = P(A\cap B) = P(A)P(B) = E[1_A]E[1_B].
\]

By linearity, (7.10) holds when $X$ and $Y$ are independent simple functions.

Suppose $X$ and $Y$ are nonnegative, measurable, and independent random variables. Let $X_n$'s and $Y_n$'s be simple nonnegative functions obtained by the method mentioned after Definition 6.3. For each $n$, $X_n$ and $Y_n$ are independent because $X_n$ is a function of $X$ and $Y_n$ is a function of $Y$, and $X$ and $Y$ are independent (Theorem 5.2). Furthermore, we have $X_nY_n \nearrow XY$ because $X_n \nearrow X$ and $Y_n \nearrow Y$. By the monotone convergence theorem,
\[
E[XY]
=
\lim_n E[X_nY_n]
=
\lim_n (E[X_n]E[Y_n])
=
(\lim_n E[X_n])(\lim_n E[Y_n]),
\]
which is the same as $E[X]E[Y]$.

Suppose $X$ and $Y$ are in $L^1(P)$. Let $X^+$ and $X^-$ denote the positive and negative parts of $X$, respectively, and let $Y^+$ and $Y^-$ denote the positive and negative parts of $Y$, respectively. We note that the pair of random variables $(X^+,X^-)$ is independent with $(Y^+,Y^-)$ because $X^+=\max(X,0)$ and $X^-=\max(-X,0)$ are functions of $X$ and $Y^+=\max(Y,0)$ and $Y^-=\max(-Y,0)$ are functions of $Y$ (Theorem 5.2). Then, we can write
\[
E[XY]=E[(X^+-X^-)(Y^+-Y^-)]
\]
\[
= E[X^+]E[Y^+] - E[X^+]E[Y^-] - E[X^-]E[Y^+] + E[X^-]E[Y^-]
\]
\[
= (E[X^+]-E[X^-])(E[Y^+]-E[Y^-])
\]
\[
= E[X]E[Y].
\]

This establishes (7.10) for nonnegative random variables $X$ and $Y$. \hfill $\square$

In view of Theorem 7.13, we can define the notion of uncorrelated random variables.

\begin{defbox}{7.3}
We say that two real-valued random variables $X$ and $Y$ are \textit{uncorrelated} if
\[
E[XY]=E[X]E[Y].
\]

The notion of uncorrelated random variables is closely related to the \textit{covariance}, which is defined as
\[
\operatorname{Cov}(X,Y)\triangleq E[(X-E[X])(Y-E[Y])].
\]

It is straightforward to show that $\operatorname{Cov}(X,Y)=0$ if and only if $X$ and $Y$ are uncorrelated. Theorem 7.13 says that two independent random variables are uncorrelated, but the converse does not hold in general.
\end{defbox}

\textbf{Example 7.5.1 (Uncorrelated But Dependent Random Variables)} \\
Let $Z$ be any discrete random variable with zero mean. Let $U$ be a random variable that equals $0$ with probability $1/2$ and $1$ with probability $1/2$. Define random variables $X$ and $Y$ by
\[
(X,Y)=
\begin{cases}
(Z,Z) & \text{if } U=0,\\
(Z,-Z) & \text{if } U=1.
\end{cases}
\]

The random variable $X$ has the same distribution as $Z$ and hence has zero mean. The expectation $E[XY]$ is zero by a symmetry argument. Therefore, $X$ and $Y$ are uncorrelated. Nevertheless, the random variables $X$ and $Y$ are not independent. It is because we can determine the absolute value of $X$ if we know the value of $Y$.
